% 星际行星（综述）
% license CCBYSA3
% type Wiki

本文根据 CC-BY-SA 协议转载翻译自维基百科\href{https://en.wikipedia.org/wiki/Rogue_planet}{相关文章}。

流浪行星（rogue planet），亦称自由漂浮行星（free-floating planet，FFP）或孤立的行星质量天体（isolated planetary-mass object，iPMO），是指一种具有行星质量、但在引力上不与任何恒星或棕矮星发生束缚关系的星际天体$^{[1][2][3][4]}$。

流浪行星可能起源于其形成所在的行星系统，并在随后被抛射出去；它们也可能在行星系统之外独立形成。仅银河系中就可能存在数十亿到数万亿颗流浪行星，这一数量范围有望由即将投入使用的南希·格雷斯·罗曼空间望远镜进一步加以精确限定$^{[5][6]}$。流浪行星进入太阳系的概率极低，更不用说对地球生命构成直接威胁——在未来一千年内发生这种情况的概率约为一万亿分之一$^{[7]}$。

一些行星质量天体可能以类似恒星的方式形成，国际天文学联合会已提议将此类天体称为亚棕矮星（sub-brown dwarfs）$^{[8]}$。一个可能的例子是 Cha~110913$-$773444，它可能是在被抛射后成为一颗流浪行星，也可能是独立形成并成为一颗亚棕矮星$^{[9]}$。
\subsection{术语}
最早的两篇发现论文分别使用了“孤立的行星质量天体”（iPMO）$^{[10]}$和“自由漂浮行星”（FFP）$^{[11]}$这两个名称。大多数天文学论文使用其中之一$^{[12][13][14]}$。“流浪行星”这一术语更常用于引力微透镜研究，而这类研究也经常使用 FFP 这一术语$^{[15][16]}$。面向公众的新闻稿则可能采用不同的名称。例如，2021 年发现至少 70 颗 FFP 时，不同新闻稿分别使用了“流浪行星”$^{[17]}$、“无恒星行星”$^{[18]}$、“游荡行星”$^{[19]}$以及“自由漂浮行星”$^{[20]}$等称呼。
\subsection{发现}
孤立的行星质量天体（isolated planetary-mass objects，iPMO）最早于 2000 年由英国研究团队 Lucas \& Roche 使用英国红外望远镜（UKIRT）在猎户座星云中发现$^{[11]}$。同年，西班牙研究团队 Zapatero Osorio 等人利用凯克望远镜的光谱观测，在 $\sigma$~猎户座星团中发现了 iPMO$^{[10]}$。对猎户座星云中这些天体的光谱研究于 2001 年发表$^{[21]}$。目前，这两个欧洲研究团队均因其准同步发现而受到认可$^{[22]}$。1999 年，日本研究团队 Oasa 等人在变色龙座 I 云中发现了一些天体$^{[23]}$，这些天体随后在 2004 年由美国研究团队 Luhman 等人通过光谱观测得到了确认$^{[24]}$。
\subsection{观测}
\begin{figure}[ht]
\centering
\includegraphics[width=6cm]{./figures/35a477b3258b7d43.png}
\caption{2021 年在上天蝎座与蛇夫座之间的区域发现了 115 个潜在的流浪行星} \label{fig_RoPl_1}
\end{figure}
发现自由漂浮行星主要有两种技术：直接成像和引力微透镜。
\subsubsection{引力微透镜}
日本大阪大学的天体物理学家角广隆宏（Takahiro Sumi）及其同事，隶属于天体物理学引力微透镜观测计划（Microlensing Observations in Astrophysics）和光学引力透镜实验（Optical Gravitational Lensing Experiment）合作组，于 2011 年发表了他们关于引力微透镜的研究成果。他们使用位于新西兰约翰山天文台的 1.8 米（5 英尺 11 英寸）MOA-II 望远镜，以及位于智利拉斯坎帕纳斯天文台的华沙大学 1.3 米（4 英尺 3 英寸）望远镜，对银河系中约 5000 万颗恒星进行了观测。他们共发现了 474 次微透镜事件，其中有 10 次事件的持续时间足够短，表明它们可能是质量约为木星质量、且在其附近没有伴随恒星的行星。研究人员根据这些观测结果估计，在银河系中，每一颗恒星大约对应近两颗木星质量的流浪行星$^{[25][26][27]}$。另有一项研究提出了一个大得多的数量估计，认为银河系中的流浪行星数量可能是恒星数量的多达 100{,}000 倍，不过该研究所涵盖的假想天体质量远小于木星$^{[28]}$。2017 年，华沙大学天文台的 Przemek~Mróz 及其同事在统计样本数量为 2011 年研究六倍的基础上开展的一项研究表明，银河系中每一颗主序星所对应的木星质量自由漂浮行星或宽轨道行星的上限为 0.25 颗$^{[29]}$。

2020 年 9 月，天文学家利用引力微透镜技术首次报告探测到一颗地球质量的流浪行星（命名为 OGLE-2016-BLG-1928），该天体不与任何恒星发生引力束缚，并在银河系中自由漂浮$^{[16][30][31]}$。
\subsubsection{直接成像}
\begin{figure}[ht]
\centering
\includegraphics[width=6cm]{./figures/250326c529d9a895.png}
\caption{使用斯皮策空间望远镜观测到的低温行星质量天体 WISE~J0830+2837（以橙色天体标示）。其温度约为 300--350~K（27--77\,$^\circ$C；80--170\,$^\circ$F）。} \label{fig_RoPl_2}
\end{figure}
微透镜行星只能通过微透镜事件本身进行研究，这使得对行星性质的刻画变得困难。因此，天文学家转而研究通过直接成像方法发现的孤立行星质量天体（iPMO）。例如，要确定一颗棕矮星或 iPMO 的质量，需要其光度和年龄等信息$^{[32]}$。然而，确定低质量天体的年龄已被证明是困难的。因此，绝大多数 iPMO 被发现于附近年轻的恒星形成区中，这些区域的年龄为天文学家所知。这些天体的年龄小于 200~Myr，质量较大（$>5\,M_J$）$^{[4]}$，并且属于 $L$ 矮星和 $T$ 矮星$^{[33][34]}$。然而，也存在一个规模虽小但正在增长的低温且年老的 $Y$ 矮星样本，其估计质量为 $8$--$20\,M_J$ $^{[35]}$。光谱型为 $Y$ 的附近流浪行星候选体包括距离为 $7.27\pm0.13$ 光年的 WISE~0855$-$0714$^{[36]}$。如果能够通过更精确的测量对这一 Y 矮星样本进行刻画，或找到更好地确定其年龄的方法，那么年老且低温的 iPMO 数量很可能会显著增加。

最早的 iPMO 于 21 世纪初通过对年轻恒星形成区的直接成像而被发现$^{[37][10][21]}$。这些通过直接成像发现的 iPMO 很可能以类似恒星的方式形成（有时被称为亚棕矮星）。也可能存在以行星方式形成、随后被抛射出去的 iPMO。然而，这类天体在运动学上应当不同于其诞生的恒星形成区，不应被环星盘所包围，并且应具有较高的金属丰度$^{[22]}$。在年轻恒星形成区中发现的 iPMO 中，没有任何一个相对于其恒星形成区表现出较高的速度。对于年老的 iPMO，低温天体 WISE~J0830+2837$^{[38]}$ 的切向速度 $V_{\mathrm{tan}}$ 约为 $100\,\mathrm{km/s}$，这一数值虽高，但仍与其在银河系中形成的情形相一致。对于 WISE~1534$-$1043$^{[39]}$，由于其约为 $200\,\mathrm{km/s}$ 的高 $V_{\mathrm{tan}}$，一种替代情景将其解释为被抛射出的系外行星，但其颜色却表明它是一颗年老、低金属丰度的棕矮星。大多数研究大质量 iPMO 的天文学家认为，它们代表了恒星形成过程的低质量端$^{[22]}$。

天文学家利用赫歇尔空间天文台和甚大望远镜观测了一颗非常年轻的自由漂浮行星质量天体 OTS~44，并证明了典型的类恒星形成模式所具有的过程同样适用于质量低至数个木星质量的孤立天体。赫歇尔的远红外观测显示，OTS~44 被一个质量至少为 10 个地球质量的盘所包围，因此最终可能形成一个大型卫星系统$^{[40]}$。利用甚大望远镜上的 SINFONI 光谱仪对 OTS~44 进行的光谱观测揭示，该盘正在积极地吸积物质，这与年轻恒星周围的盘相似$^{[40]}$。
\subsubsection{双星}
\begin{figure}[ht]
\centering
\includegraphics[width=6cm]{./figures/2fb0b35b79287356.png}
\caption{2MASS~J1119–1137AB，即首个被发现的行星质量双星系统，位于 TW~Hydrae 协会中} \label{fig_RoPl_3}
\end{figure}
\begin{figure}[ht]
\centering
\includegraphics[width=6cm]{./figures/4a9c799a9b362b83.png}
\caption{JuMBO~29，一个候选的 $12.5+3\,M_J$ 行星质量双星系统，分离距离为 135~AU，位于猎户座星云中} \label{fig_RoPl_4}
\end{figure}
首个被分辨开的行星质量双星的发现对象是 2MASS~J1119–1137AB。然而，目前还已知其他一些双星系统，例如 2MASS~J1553022+153236AB$^{[41][42]}$、WISE~1828+2650、WISE~0146+4234、WISE~J0336–0143（也可能是一个由棕矮星与行星质量天体组成的双星系统（BD+PMO））、NIRISS-NGC1333-12$^{[43]}$，以及 Zhang 等人发现的若干天体$^{[42]}$。

在猎户座星云中，人们发现了由 40 个宽双星系统和 2 个三体系统构成的一群天体。这一发现之所以令人惊讶，有两个原因：其一，针对棕矮星双星的理论趋势预测，随着质量降低，低质量天体之间的分离距离应当减小；其二，还预测双星比例会随质量降低而减少。这些双星被命名为木星质量双星天体（Jupiter-mass Binary Objects，JuMBOs）；它们至少占所有 iPMO 的 9\%，其分离距离小于 340~AU$^{[44]}$。目前尚不清楚这些 JuMBO 的形成机制，但一项深入研究认为它们是在原地形成的，方式类似于恒星$^{[45]}$。如果它们确实以恒星的方式形成，那么必然存在某种尚未明确的“额外因素”使其得以形成；如果它们以行星的方式形成并随后被抛射出去，则需要解释为何这些双星在抛射过程中并未解体。未来利用 JWST 的观测或许能够判定这些天体究竟是以被抛射行星的方式形成，还是以恒星的方式形成$^{[44]}$。Kevin~Luhman 对 NIRCam 数据进行了重新分析，发现大多数 JuMBO 并未出现在他所选取的亚恒星天体样本中。此外，其颜色与被红化的背景源或低信噪比源一致。他仅认为 JuMBO~29 是一个可靠的行星质量双星系统候选体$^{[46]}$。
\subsubsection{已知 iPMO 的总数}
目前已知的候选 iPMO 很可能有数百个$^{[47][44]}$，其中有一百多个天体已获得光谱观测结果$^{[48][49][50]}$，此外还有数量虽少但正在增加的通过引力微透镜发现的候选体。一些大型巡天项目包括：

截至 2021 年 12 月，迄今为止规模最大的流浪行星样本被发现，其数量至少为 70 个，在假定不同年龄的情况下最多可达 170 个。这些天体位于上天蝎座与蛇夫座之间的 OB 协会中，质量介于$4$ 至 $13\,M_J$之间，年龄约为 300 万至 1000 万年，最有可能通过气体云的引力坍缩形成，或在原行星盘中形成后因动力学不稳定而被抛射出去$^{[47][17][51][19]}$。随后利用昴星团望远镜和加那利大型望远镜开展的光谱跟进观测表明，该样本的污染率相当低（$\leq6\%$）。其中 16 个年轻天体的质量介于$3$ 至 $14\,M_J$ 之间，确认它们确实属于行星质量天体$^{[50]}$。

2023 年 10 月，利用 JWST 在梯形星团及猎户座星云内部区域又发现了一个更大的行星质量天体候选体样本，总数达到 540 个。这些天体的质量范围为 $13$ 至 $0.6\,M_J$。其中相当一部分天体形成了宽双星系统，这一现象此前并未被预测到$^{[44]}$。
\subsection{形成}
总体而言，孤立行星质量天体（isolated planetary-mass object，iPMO）的形成主要有两种情景：其一，它像行星一样在恒星周围形成，随后被抛射出去；其二，它像低质量恒星或棕矮星一样在孤立环境中形成。这两种形成方式会影响其成分和运动特性$^{[22]}$。

近期研究表明，流浪行星既可能通过恒星形成区内的直接引力坍缩形成，也可能在其诞生的行星系统中形成后被抛射出来，并在随后与已建立的行星系统发生相互作用，从而影响这些系统的轨道结构和整体统计特征。许多此类天体很可能最初起源于行星系统之中，随后在动力学过程中被驱逐，而另一些则可能是在孤立环境中形成的。除了在近距离遭遇或潜在的俘获事件中改变系统稳定性之外，流浪行星还可能输送挥发物，从而增强前生物化学过程，并创造有利于提高生物多样性的条件。这些形成机制、动力学过程、生化作用以及生态效应的综合影响，在塑造系外行星系统的分布与演化方面发挥着重要作用$^{[52]}$。
\subsubsection{类似恒星的形成}
2001 年的模型认为，质量至少为一个木星质量的天体可以通过分子云的坍缩与碎裂过程形成$^{[53]}$。JWST 之前的观测表明，质量低于$3$--$5\,M_J$ 的天体不太可能独立形成$^{[4]}$。2023 年，利用 JWST 对梯形星团的观测显示，质量低至 $0.6\,M_J$ 的天体也可能独立形成，并不需要一个陡峭的质量截断$^{[44]}$。一种称为小球状体（globulettes）的特殊云核被认为是棕矮星和行星质量天体的诞生地。小球状体可见于玫瑰星云和 IC~1805 中$^{[54]}$。有时，年轻的 iPMO 仍被一个可能形成系外卫星的盘所包围。由于这种系外卫星围绕其宿主行星的轨道非常紧密，其发生凌星现象的概率高达$10$--$15\%$，$^{[55]}$。

\textbf{盘}

一些非常年轻的恒星形成区（通常年龄小于 500 万年）有时包含具有红外过量并显示吸积迹象的孤立行星质量天体。其中最著名的是位于变色龙座 I 的 iPMO——OTS~44，其被发现具有一个盘。变色龙座 I 和 II 中还存在其他带盘的 iPMO 候选体$^{[56][57][33]}$。其他包含带盘或吸积 iPMO 的恒星形成区包括：天狼座 I$^{[57]}$、蛇夫座 $\rho$ 云复合体$^{[58]}$、$\sigma$~猎户座星团$^{[59]}$、猎户座星云$^{[60]}$、金牛座$^{[58][61]}$、NGC~1333$^{[62]}$ 以及 IC~348$^{[63]}$。利用 ALMA 对棕矮星和 iPMO 周围盘的大型巡天发现，这些盘的质量不足以形成地球质量的行星，但仍存在这些盘中已形成行星的可能性$^{[58]}$。对红矮星的研究表明，其中一些在相对较老的年龄阶段仍保留富含气体的盘，这类盘被称为“彼得·潘盘”（Peter Pan Disks），而这一趋势可能延伸至行星质量范围。一个彼得·潘盘的例子是年龄约为 45~Myr、质量约为$13.7\,M_J$ 的棕矮星 2MASS~J02265658-5327032，其质量已接近行星质量范围$^{[64]}$。近期对附近行星质量天体 2MASS~J11151597+1937266 的研究发现，该 iPMO 周围存在一个盘，并显示出来自盘的吸积迹象以及红外过量$^{[65]}$。2025 年 5 月，研究人员利用 JWST 发现 Cha~1107-7626 周围的盘中含有烃类物质。Cha~1107-7626（$6--10\,M_J$）是已知具有尘埃盘的最低质量天体之一$^{[66]}$。进一步的 JWST 光谱观测表明，硅酸盐和烃类是行星质量天体盘中的常见成分。这些盘显示出明显的尘粒生长和结晶化证据，与棕矮星和恒星周围盘中所见的现象相似，表明这些盘具备形成岩质伴体的能力$^{[67]}$。
\subsubsection{类似行星的形成}
被抛射行星在理论上预计主要为低质量天体（$<30\,M_\oplus$，见 Ma 等人的图 1）$^{[68]}$，其平均质量取决于其宿主恒星的质量。Ma 等人$^{[68]}$ 的数值模拟表明，质量为 $1\,M_\odot$ 的恒星中，有 $17.5\%$ 会发生行星抛射事件，每颗恒星平均抛射出总质量为 $16.8\,M_\oplus$ 的物质，而单个自由漂浮行星（FFP）的典型（中位）质量约为 $0.8\,M_\oplus$。对于质量为 $0.3\,M_\odot$ 的低质量红矮星，有 $12\%$ 的恒星会发生抛射事件，每颗恒星平均抛射出总质量为 $5.1\,M_\oplus$ 的物质，而单个 FFP 的典型质量约为 $0.3\,M_\oplus$。

Hong 等人$^{[69]}$ 预测，系外卫星可以在行星—行星相互作用中被散射并成为被抛射的系外卫星。质量较高（$0.3$--$1\,M_J$）的被抛射 FFP 在理论上也是可能的，但预计十分罕见$^{[68]}$。行星的抛射可以通过行星—行星散射过程发生，或由恒星近距离掠过所触发。另一种可能性是原行星盘碎片被抛射出去，并随后演化形成行星质量天体$^{[70]}$。还有一种提出的情景是，位于倾斜的环双星轨道上的行星被抛射出去。中央双星与行星之间以及行星彼此之间的相互作用，可能导致系统中质量较低的行星被抛射$^{[71][72]}$。然而，这一机制的有效性依赖于相遇的几何构型，而该构型在观测和理论上目前均缺乏良好约束。
\subsubsection{通过年轻环星盘相互作用形成}
处于边缘引力稳定状态的年轻环星盘之间的相互作用，可能产生拉长的潮汐桥结构，这些结构可在局部发生坍缩并形成 iPMO$^{[73]}$。这些 iPMO 拥有与观测结果相符的大尺度盘结构$^{[60]}$，而被抛射行星假说难以解释这一特征。正如在梯形星团中的 iPMO 所显示的那样，它们在形成时还具有较高的多体系统比例$^{[44]}$。然而，这一机制的有效性同样依赖于相遇几何构型，而该构型在观测和理论上目前仍缺乏良好约束$^{[74]}$。
\subsubsection{其他情景}
如果恒星或棕矮星胚胎的吸积过程被中途终止，其质量可能保持在足够低的水平，从而成为行星质量天体。这种吸积中断可能发生在胚胎被抛射的情况下，或其环星盘在$O$型恒星附近经历光致蒸发时。通过“被抛射胚胎”情景形成的天体，其盘结构会较小甚至不存在，并且这类天体的双星比例会降低。也有可能，自由漂浮的行星质量天体是多种形成机制共同作用的结果$^{[70]}$。
\subsection{命运}
大多数孤立的行星质量天体将永远在星际空间中漂浮。

部分 iPMO 会与行星系统发生近距离相遇。这种罕见的相遇可能产生三种结果：iPMO 仍保持不受束缚状态；它可能与恒星形成弱引力束缚；或者它可能将一颗系外行星“踢出”系统并取而代之。数值模拟表明，绝大多数此类相遇都会导致一次俘获事件，其中 iPMO 与恒星形成弱束缚状态，其引力束缚能较低，并处在拉长的高偏心率轨道上。这类轨道并不稳定，其中约 90\% 的天体会因行星—行星相互作用而获得能量，最终再次被抛射回星际空间。只有约 1\% 的恒星会经历这种短暂的俘获过程$^{[75]}$。
\subsection{热量}
\begin{figure}[ht]
\centering
\includegraphics[width=6cm]{./figures/e7e727d8b9ebfbed.png}
\caption{一颗木星大小流浪行星的艺术构想图} \label{fig_RoPl_5}
\end{figure}
星际行星几乎不产生热量，也不会受到恒星的加热$^{[76]}$。然而，1998 年 David~J.~Stevenson 提出理论认为，一些在星际空间中漂流的行星尺度天体可能维持一层厚重的大气而不会冻结。他提出，这类大气能够依靠富含氢的大气层在高压条件下产生的远红外辐射不透明性而得以保存$^{[77]}$。

在行星系统形成过程中，若干较小的原行星天体可能会被抛射出系统$^{[78]}$。被抛射的天体将接收到更少由恒星产生的紫外辐射，而这些紫外辐射通常会剥离大气中的轻元素。即便是地球大小的天体，其引力也足以阻止大气中氢和氦的逃逸$^{[77]}$。对于地球尺度的天体而言，其核心中残余放射性同位素衰变所产生的地热能，可能使其表面温度维持在水的熔点以上$^{[77]}$，从而允许液态水海洋的存在。这类行星很可能在很长时间尺度上保持地质活动性。如果它们具有由地磁发电机产生的保护性磁层以及海底火山活动，那么热液喷口就可能为生命提供能量来源$^{[77]}$。由于其热微波辐射极其微弱，这些天体将难以被探测；不过，对于距离地球小于 1000 个天文单位的天体，其反射的太阳辐射以及远红外热辐射仍有可能被探测到$^{[79]}$。约有 $5\%$ 的地球大小被抛射行星在抛射后仍能保留其月球尺度的天然卫星。一个较大的卫星将成为显著的地质潮汐加热来源$^{[80]}$。
\subsection{列表}
下表列出了已确认或疑似的流浪行星。这些行星究竟是从围绕恒星的轨道中被抛射出去的，还是作为亚棕矮星而独立形成，目前尚不清楚。对于质量异常低的流浪行星（例如 OGLE-2012-BLG-1323 和 KMT-2019-BLG-2073），它们是否具备独立形成的能力，目前仍未知。
\subsubsection{通过直接成像发现}
这些天体是通过直接成像方法发现的。其中许多是在年轻的恒星团或恒星协会中被发现的，也已知存在少数年老的天体（例如 WISE~0855-0714）。该列表按发现年份排序。

\begin{table}[ht]
\centering
\caption{请输入表格标题}\label{tab_RoPl1}
\begin{tabular}{|c|c|c|c|c|c|c|c|}
\hline
\textbf{系外行星} & \textbf{质量（$M_J$）}& \textbf{年龄（Myr）} &\textbf{ 距离（ly）} & \textbf{光谱型} & \textbf{状态} & \textbf{恒星协会成员关系} & \textbf{发现}\\
\hline
OTS~44 & $\sim11.5$ & $0.5-3 $& 554 & M9.5 & 很可能是低质量棕矮星$^{[37]}$ & 变色龙座~I & 1998\\
\hline
S~Ori~52 & $2-8$ & $1-5 $& $1{,}150$  &  & 年龄和质量不确定；可能是前景棕矮星 & $\sigma$~猎户座星团 & 2000$^{[10]}$\\
\hline
Proplyd~061-401 & $\sim11$ & 1 & $1{,}344$ & $L4-L5$ & 候选体；本研究共发现 15 个候选体 & 猎户座星云 & 2001$^{[21]}$\\
\hline
S~Ori~70 & 3 & 3 & 1150 & T6 & 可能为混入天体？$^{[22]}$ & $\sigma$~猎户座星团 & 2002\\
\hline
Cha~110913-773444 & $5-15$ & $\sim2$ & 529 & $>$M9.5 & 已确认 & 变色龙座~I & 2004$^{[81]}$\\
\hline
SIMP~J013656.5+093347 & $11-13$ & $\sim200$ & $20-22$ & T2.5 & 候选体 & 船底座邻近运动群 & 2006$^{[82][83]}$\\
\hline
Cha~1107-7626 & $6-10$ & $1-5$ & $620$ & $L0-L1$ & 已确认 & 变色龙座~I & 2008$^{[84]}$\\
\hline
UGPS~J072227.51-054031.2 & $0.66-16.02$,$^{[85][86]}$ & 1000--5000& 13 & T9 & 质量不确定 & 无 & 2010\\
\hline
M10$-$4450 & $2-3$ & 1 & 325 & T & 候选体 & $\rho$~蛇夫座分子云 & 2010$^{[87]}$\\
\hline
WISE~1828+2650 & $3-6$ 或 $0.5-20$,$^{[88]}$ & $2-4$ 或 $0.1-10$,$^{[88]}$ & 47 & >Y2 & 候选体，可能为双星 & 无 & 2011\\
\hline
CFBDSIR~2149-0403 & $4-7$ & $110-130$ & $117-143$ & T7 & 候选体 & AB~Doradus 运动群 & 2012$^{[89]}$\\
\hline
SONYC-NGC1333-36 & $\sim6$ & 1 & 978 & L3 & 候选体；NGC~1333 中另有两个质量低于 $15\,M_J$ 的天体 & NGC~1333 & 2012$^{[90]}$\\
\hline
SSTc2d~J183037.2+011837 & $2-4$ & 3 & $848-1354$ & T? & 候选体，亦称 ID~4  & 巨蛇座核心星团$^{[91]}$（位于巨蛇座分子云） & 2012$^{[12]}$\\
\hline
PSO~J318.5-22 & $6.24-7.60$,$^{[85][86]}$ & $21-27$ & 72.32 & L7 & 已确认；亦称 2MASS~J21140802-2251358 & $\beta$~绘架座运动群 & 2013$^{[14][92]}$\\
\hline
2MASS~J2208+2921 & $11-13$  & $21-27$ & 115 & L3$\gamma$ & 候选体；需径向速度确认 & $\beta$~绘架座运动群 & 2014$^{[93]}$\\
\hline
WISE~J1741-4642 & $4-21$ & $23-130$ &  & L7pec  & 候选体  & $\beta$~绘架座或 AB~Doradus 运动群 & 2014$^{[94]}$\\
\hline
WISE~0855-0714 & $3-10$ & $>1{,}000$ & 7.1 & Y4 & 年龄不确定，但由于位于太阳邻域，推测为年老天体；$^{[95]}$ 即便在 $12$~Gyr 的高龄情形下仍为候选体（宇宙年龄为 $13.8$~Gyr）。目前已知距离最近的可能流浪行星 & 无 & 2014$^{[96]}$\\
\hline
2MASS~J12074836-3900043 & $\sim15$,$^{[97]}$ & 7--13 & 200 & L1 & 候选体；需更精确的距离测量 & TW~Hydrae 协会$^{[98]}$ & 2014$^{[99]}$\\
\hline
SIMP~J2154-1055 & 9--11 & 30--50 & 63 & L4$\beta$ & 年龄存在争议$^{[100]}$ & Argus 协会 & 2014$^{[101]}$\\
\hline
SDSS~J111010.01+011613.1 & $10.83-11.73$,$^{[85][86]}$ & $110-130$ & 63 & T5.5 & 已确认$^{[85]}$ & AB~Doradus 运动群 & 2015$^{[34]}$\\
\hline
2MASS~J11193254-1137466~AB & $4-8$ & $7-13$ & $\sim90$ & L7 & 双星候选体，其中一个分量具有候选系外卫星或可变大气$^{[55]}$ & TW~Hydrae 协会 & 2016$^{[102]}$\\
\hline
WISEA~1147 & $5-3 $& $7-13$ & $\sim100$ & L7 & 候选体 & TW~Hydrae 协会  & 2016$^{[13]}$\\
\hline
USco~J155150.2-213457 & $8-10$ & $6.907-10$ & 104 & L6 & 候选体，低表面引力 &上天蝎座协会 & 2016$^{[103]}$\\
\hline
Proplyd~133-353 & $<13$ & $0.5-1$ & $1{,}344$ & M9.5 & 具有光致蒸发盘的候选体 & 猎户座星云 & 2016$^{[60]}$\\
\hline
Cha~J11110675-7636030 & $3-6$ & $1-3$ & $520-550$  & M9-L2 & 候选体，但可能被盘包围，从而可能为亚棕矮星；该研究还提出了其他候选体 & 变色龙座~I & 2017$^{[33]}$\\
\hline
PSO~J077.1+24 & 6 & $1-2$ & 470 & L2 & 候选体；该研究还在金牛座中发表了另一候选体 & 金牛座分子云 & 2017$^{[104]}$\\
\hline
2MASS~J1115+1937 & $6^{+8}_{-4}$ & $5-45$ & 147  & L2$\gamma$ & 具有吸积盘  & 场星，可能为被抛射天体 & 2017\\
\hline
Calar~25 & $11-12$ & 120 & 435  &  & 已确认 & 昴宿星团 & 2018$^{[105]}$\\ 
\hline
2MASS~J1324+6358 & $10.7-11.8$ & $\sim150$ & $\sim33$ & T2 & 异常偏红且不太可能为双星；可靠候选体$^{[85][86]}$ & AB~Doradus 运动群 & 2007,~2018$^{[106]}$\\
\hline
WISE~J0830+2837 & $4-13$ & $>1{,}000$ & $31.3-42.7$ & $>$Y1 & 年龄不确定，但由于速度较高（高 $V_{\mathrm{tan}}$ 指示年老恒星族群）而被认为是年老天体；若年龄小于 10~Gyr 则为候选体 & 无 & 2020$^{[38]}$\\
\hline
2MASS~J0718-6415 & $3\pm1$ & $16-28$ & 30.5 & T5 & BPMG 候选成员；自转周期极短，为 $1.08$ 小时，与棕矮星 2MASS~J0348-6022 相当$^{[107][108]}$ & $\beta$~绘架座运动群 & 2021\\
\hline
DANCe~J16081299-2304316 & $3.1-6.3$ & $3-10$ & 104 & L6 & 该研究发表的至少 70 个候选体之一，光谱与 HR~8799c 相似 & 上天蝎座协会 & 2021$^{[47][50]}$\\
\hline
WISE~J2255-3118 & $2.15-2.59$ & 24 & $\sim45$ & T8 & 非常偏红，候选体$^{[85][86]}$，已确认？$^{[109]}$ & $\beta$~绘架座运动群 & 2011,~2021$^{[49]}$\\
\hline
WISE~J024124.73-365328.0 & $4.64-5.30$ & 45 & $\sim61$ & T7 & 候选体$^{[85][86]}$ & Argus 协会 & 2012,~2021$^{[49]}$\\
\hline
2MASS~J0013-1143 & $7.29-8.25$ & 45 & $\sim82$ & T4 & 双星候选体或复合大气；候选体$^{[85][86]}$ & Argus 协会 & 2017,~2021$^{[49]}$\\
\hline
SDSS~J020742.48+000056.2 & $7.11-8.61$ & 45 & $\sim112$ & T4.5 & 候选体$^{[85][86]}$ & Argus 协会 & 2002,~2021$^{[49]}$\\
\hline
2MASSI~J0453264-175154 & $12.68-12.98$ & 24 & $\sim99$ & L2.5$\beta$ & 低表面引力，候选体$^{[85][86]}$ & $\beta$~绘架座运动群 & 2003,~2023$^{[85][86]}$\\
\hline
CWISE~J0506+0738 & $7\pm2$ & 22 & 104 & $L8\gamma-T0\gamma$ & BPMG 候选成员；近红外颜色极端偏红$^{[110]}$ & $\beta$~绘架座运动群 & 2023\\
\hline
\end{tabular}
\end{table}
\subsubsection{通过引力微透镜发现}
这些天体是通过引力微透镜方法发现的。通过引力微透镜发现的流浪行星只能依赖透镜事件本身进行研究。其中一些天体也可能是在一颗未被观测到的恒星周围处于宽轨道上的系外行星$^{[111]}$。
\begin{table}[ht]
\centering
\caption{请输入表格标题}\label{tab_RoPl2}
\begin{tabular}{|c|c|c|c|c|c|}
\hline
\textbf{系外行星} & \textbf{质量（$M_J$}） & \textbf{质量（$M_\oplus$）}  & \textbf{距离（ly}） & \textbf{状态}  & \textbf{发现} \\
\hline
KMT-2023-BLG-2669 & $0.025-0.25$ & $8-80$ & & 候选体；需确定距离 & 2024$^{[112]}$\\
\hline
OGLE-2012-BLG-1323 & $0.0072-0.072$ & $2.3-23$ & &  候选体；需确定距离 & 2017$^{[113][114][115][116]}$\\
\hline
OGLE-2017-BLG-0560 & $1.9-20$ & $604-3{,}256$ &  & 候选体；需确定距离 & 2017$^{[114][115][116]}$\\
\hline
MOA-2015-BLG-337L & 9.85 & $3{,}130$ & $23{,}156$ & 可能是一个棕矮星双星系统 & 2018$^{[117][118]}$\\
\hline
KMT-2019-BLG-2073 & 0.19 & 59 &  & 候选体；需确定距离 & 2020$^{[119]}$\\
\hline
OGLE-2016-BLG-1928 & $0.001-0.006$ & $0.3-2$ & $30{,}000-180{,}000$  & 候选体 & 2020$^{[111]}$\\
\hline
OGLE-2019-BLG-0551 & $0.0242-0.3$ & $7.69-95$ &  &  刻画程度较低$^{[120]}$ & 2020$^{[120]}$\\
\hline
VVV-2012-BLG-0472L & 10.5 & $3{,}337$ & $3{,}200$ &  & 2022$^{[121]}$\\
\hline
MOA-9y-770L & 0.07 & $22.3^{+42.2}_{-17.4}$ & $22{,}700$ & & 2023$^{[122]}$\\
\hline
MOA-9y-5919L & 0.0012 或 0.0024 & $0.37^{+1.11}_{-0.27}$ 或 $0.75^{+1.23}_{-0.46}$ & $14{,}700$ 或 $19{,}300$ &  & 2023$^{[122]}$\\
\hline
OGLE-2017-BLG-1170L & $3.06^{+1.34}_{-1.16}$ &  & $24{,}700$ & 候选体 & 2019$^{[123]}$\\
\hline
OGLE-2017-BLG-1170L & $1.85^{+0.79}_{-0.70}$ & &  $24{,}700$ & 候选体  & 2019$^{[123]}$\\
\hline
\end{tabular}
\end{table}




